\subsection{DRG.3.3 en CTR: estados distinguibles en el modelo de cifrado ideal}
\label{sec:ctr-effective-struct}
Fijamos un contador inicial \(v\), \(M=2^{128}\) y \(g=2^{246}\).
Dos estados con idéntico comportamiento ante todas las entradas permitidas
deben producir el mismo primer bloque para cada una de cinco peticiones
alternativas \(A=0,1,2,3,4\). Compararemos las cinco respuestas que
\emph{daría} el estado inicial, sin hacer cinco llamadas sucesivas desde
estados distintos ni conceder acceso de restauración al observador.
El objetivo estructural es encontrar \(g\) estados con firmas distintas.
La contabilidad siguiente toma AES como una permutación aleatoria
independiente por clave. Es una afirmación de cardinalidad del perfil
algorítmico \cite[§§91--103, 115, 207 y 1007]{ais}.

\paragraph{Selección de los representantes.}
Las \(2^{255}\) claves antiguas con bit alto cero se reparten en \(g\)
grupos disjuntos de 512. En esta fase sólo se consultan tablas de claves
antiguas. Para cada candidata, las dos primeras salidas de
\(G(k,v)\), que constituyen la clave efectiva de la petición \(k'\),
son un par uniforme entre los \(H=M(M-1)\) pares distintos.
Exactamente la mitad de esos pares tiene bit alto uno.
De cada grupo elegimos la primera candidata cuya \(k'\) tenga bit alto
uno y sea distinta de las elegidas en grupos anteriores.
En cualquier historial de la selección hay a lo sumo \(g\) claves
prohibidas. Cada candidata se prueba una sola vez bajo una clave antigua
nueva. Su probabilidad condicional de fallo queda acotada por
\[
   \frac{H/2+g}{H}\le\frac9{16}.
\] El tercer bloque de \(G(k,v)\) da un desplazamiento
del contador que puede fijarse junto con las tablas de esta fase.
La probabilidad de que algún grupo carezca de representante es a lo sumo
\(g(9/16)^{512}\). La cota usa productos de probabilidades condicionales
dentro de un grupo y una unión sobre grupos.
\proofat{EffectiveCTR.scaledGroupFailure}{scaledGroupFailure} formaliza
el conteo de rutas fallidas con historial arbitrario.

\paragraph{Separación por las cinco respuestas.}
Condicionadas a las tablas de la primera fase, las \(g\) claves derivadas
elegidas son altas y diferentes. Sus tablas siguen siendo permutaciones
uniformes e independientes, pues ninguna se ha consultado. Para cada una,
los cinco argumentos de AES son \(\operatorname{inc}(V'\oplus A)\),
con \(V'=G_V(k,v)\) ya fijado. Son cinco puntos distintos.
\proofat{EffectiveCTR.signatureUnderDerivedKey}{signatureUnderDerivedKey}
y \proofat{EffectiveCTR.pointInjective}{pointInjective} vinculan el argumento
a \texttt{Generate} de 90A. Una segunda firma coincide con una primera
fijada con probabilidad \(1/(M)_5\), donde \((M)_5=M(M-1)\cdots(M-4)\). Quedan \((M-5)!\) permutaciones que realizan las cinco asignaciones,
entre \(M!\) posibles. La probabilidad total de fallo cumple
\[
    g(9/16)^{512}+\frac{g(g-1)}{2(M)_5}<2^{-148}.
\]
\proofat{EffectiveCTR.totalFailure}{totalFailure} verifica la desigualdad por
aritmética entera exacta.
\proofat{EffectiveCTR.separateClasses}{separateClasses} muestra que las firmas
distintas implican clases de comportamiento distintas. Cada par candidato
\((k,v)\) es una imagen de la siembra sin DF si se admite todo el alfabeto
de material de 384 bits
(\proofat{CTRProfile.seed_attainable}{seed\_attainable}).

Este resultado obtiene \(2^{246}\) clases distinguibles mediante entradas
adicionales permitidas, con la probabilidad indicada sobre la familia
ideal. Para la instancia deben justificarse los estados que permite su
fuente real, la aplicación del modelo a AES y la dificultad de inferir
el estado efectivo. Las cinco peticiones son experimentos alternativos
sobre la semántica del perfil. Una interfaz restringida permanentemente
a la entrada vacía necesitaría una separación de estados con esa restricción.
La predicción autónoma sigue el análisis computacional del expediente CTR.
